\section{Introduction}\label{sec:introduction}

Wall's finiteness conditions ask when algebraic restrictions on a CW
complex force it to have a homotopy model of a prescribed dimension.
Let $X$ be a finite connected CW complex, let $G=\pi_1(X)$, and let
$\widetilde X$ be its universal cover. The $D(2)$ condition is
\begin{equation}\label{eq:d2-condition}
 H_i(\widetilde X;\Z)=0\quad(i>2),\qquad
 H^3(X;M)=0\quad\text{for every finitely generated }\Z G\text{-module }M,
\end{equation}
where cohomology uses the corresponding local coefficients.
The finite $D(2)$ problem asks whether every such $X$ is homotopy
equivalent to a finite CW complex of dimension at most two
\cite[Section~2]{Wall65}, \cite[Problem~D3]{Wall79}.
The finiteness of both complexes is part of the question, and no
asphericity hypothesis is imposed. This is also the formulation in
Hofmann--Nicholson~\cite[p.~254]{HofmannNicholson27}.

\begin{theorem}\label{thm:main}
There is a finite connected three-dimensional CW complex $X$ such that
\[
 H_i(\widetilde X;\Z)=0\quad(i>2),\qquad
 H^3(X;M)=0\quad\text{for every }\Z[\pi_1(X)]\text{-module }M,
\]
but $X$ is not homotopy equivalent to any finite CW complex of
dimension at most two.
\end{theorem}

The example has infinite fundamental group; the finiteness assertion
concerns its cell structure.

\subsection{Background and approach}

Wall introduced the dimension conditions in his work on finiteness
for CW complexes~\cite{Wall65}. His reduction produces a finite
three-dimensional model in the $D(2)$ case; the remaining question is
whether its three-dimensional cells can also be removed up to homotopy.
The problem reappeared as Problem~D3 in his 1979 list~\cite{Wall79}.
It lies at the interface between two-dimensional topology and the
realization of chain complexes over group rings. Johnson's work
\cite{Johnson03} develops this algebraic viewpoint and the role of
cancellation in passing from a stable realization to an actual
two-dimensional complex.

A stable form of the problem has a positive answer: adjoining
sufficiently many two-spheres makes a finite $D(2)$ complex homotopy
equivalent to a finite two-dimensional complex. This goes back to
Cohen~\cite{Cohen78}; Hambleton~\cite[Lemma~2.1]{Hambleton18} gives
a proof with the number of added spheres bounded by the number of
three-dimensional cells. Such stabilization does not answer the
unstabilized question. There are also substantial positive results
for particular fundamental groups. Hambleton~\cite[Theorem~B]{Hambleton18}
proves the $D(2)$ property for finite subgroups of $\operatorname{SO}(3)$.
More recently, Hofmann and Nicholson~\cite[Theorem~A]{HofmannNicholson27}
establish the $D(2)$ property for the generalized quaternion groups of
orders $24$ and $32$ and recover the order-$28$ case previously proved
using work of Mannan--Popiel and Nicholson~\cite{MannanPopiel21,Nicholson21}.
In particular,
they show that the longstanding candidate
with fundamental group of order $32$, proposed by Cohen and Dyer,
does have a finite two-dimensional model.

Candidate counterexamples have also come from relation modules.
A relation gap occurs when the relator subgroup of a free group needs
more normal generators than its abelianization needs as a module over
the quotient group ring.
Bridson and Tweedale~\cite[Sections~3--4]{BridsonTweedale07} use
coprime annihilation relations to reduce the number of module generators
and construct finite $D(2)$ complexes with virtually free fundamental
groups. Their failure to have finite two-dimensional models is conditional
on a relation-gap assertion.

We apply Quillen's plus construction~\cite{Quillen71}. It kills a
perfect normal subgroup of the fundamental
group while preserving homology with coefficients from the quotient.
Mannan~\cite[Section~2]{Mannan09} shows directly that applying it to a
finite two-dimensional complex along a finitely normally generated
perfect subgroup produces a finite $D(2)$ complex. He also proves
that every finite cohomologically two-dimensional three-dimensional
complex arises in this way up to homotopy
\cite[Theorem~3.4]{Mannan09}. Thus this construction supplies the
dimension condition; the task is to find an obstruction to a
two-dimensional model of its result.

Our obstruction uses a character. For a connected space $Y$, a
homomorphism $\rho:\pi_1(Y)\to\C^\times$ determines a
one-dimensional complex local system $L_\rho$. We prove that, when
$Y$ is a finite two-dimensional complex,
\begin{equation}\label{eq:obstruction-overview}
 H_2(Y;L_\rho)=0
 \quad\Longrightarrow\quad
 \rho(g)=1\ \text{for every finite-order }g\in\pi_1(Y).
\end{equation}
The proof gives the stronger statement that every map $F\to Y$ from
a finite connected simplicial complex with $b_1(F;\Q)=0$ pulls $L_\rho$
back to a system with trivial monodromy. The method is a tower of
infinite cyclic covers and image subcomplexes, following the cyclic-tower
method of Howie~\cite{Howie81Pairs}; see also the formulation of
Louder--Wilton~\cite[Section~3]{LouderWilton17} and its account of
the earlier work of Papakyriakopoulos and Stallings. The additional input
is that vanishing of rank-one twisted second homology persists under
these cyclic covers, by a Laurent-polynomial rank argument.

\subsection{Proof structure}

Section~\ref{sec:tower} establishes this obstruction independently
of the construction. If a map $F\to Y$ violated it, its finite image
would have nonpositive Euler characteristic and hence admit an
infinite cyclic cover. The map lifts because $b_1(F;\Q)=0$.
Twisted second homology still vanishes upstairs, so the argument
can be repeated. Each lifted image contains strictly more simplices,
although all are images of the same fixed triangulated $F$. This
forces a contradiction.

Section~\ref{sec:plus} presents the finite plus construction with
its explicit cellular splitting. Section~\ref{sec:construction}
then gives a five-generator, four-relator presentation and a perfect
normal subgroup generated normally by two powers. The quotient has
an element of order two detected by a complex character. An exact
four-row boundary calculation shows that the presentation complex
has zero second homology with that character. The plus construction
preserves this vanishing and gives the required finite $D(2)$
complex. Equation~\eqref{eq:obstruction-overview} excludes its
two-dimensional realization.

The geometric obstruction and the finite plus construction play separate
roles: the former applies to any rank-one local system on a finite
two-dimensional complex, while the latter is standard once a suitable
perfect normal subgroup is available. Their combination makes the
failure of two-dimensional realization visible in a small twisted
cellular boundary matrix.
